\chapter{Fondements Théoriques}
\label{chap:theorie}

\section{Pare-feu FortiGate}

\subsection{Définition et évolution}

Un pare-feu FortiGate est un dispositif de sécurité réseau qui combine les fonctionnalités traditionnelles des pare-feux (packet filtering, NAT) avec des capacités avancées d'inspection et de contrôle du trafic \cite{1}.

Contrairement aux pare-feux de première génération qui se basent uniquement sur les adresses IP et ports, les pare-feux FortiGate offrent :
\begin{itemize}
    \item Deep Packet Inspection (DPI) des paquets
    \item L'identification des applications au-delà des ports et protocoles
    \item L'intégration de modules UTM (Unified Threat Management)
    \item Le support du routage dynamique (OSPF, BGP)
    \item High Availability (HA)
\end{itemize}

\subsection{Architecture FortiGate}

Les pare-feux FortiGate de Fortinet utilisent une architecture matérielle spécialisée basée sur des processeurs dédiés :

\begin{figure}[H]
\centering
\begin{tikzpicture}[
    node distance=1cm,
    block/.style={rectangle, draw, fill=blue!10, text width=2.5cm, minimum height=0.9cm, align=center},
    spu/.style={rectangle, draw, fill=red!10, text width=2.2cm, minimum height=0.7cm, align=center},
    scale=0.8,
    transform shape
]
    % CPU principal
    \node[block] (cpu) {CPU Principal\\(x86)};
    
    % NP7
    \node[spu, above left=of cpu] (np7) {NP7\\(Network Processor)};
    
    % CP9
    \node[spu, above right=of cpu] (cp9) {CP9\\(Content Processor)};
    
    % Memory
    \node[block, below=of cpu] (mem) {Mémoire\\(RAM)};
    
    % Network interfaces
    \node[block, left=of np7] (net1) {Interfaces\\Réseau};
    \node[block, right=of cp9] (net2) {Interfaces\\Réseau};
    
    % Arrows
    \draw[<->, thick] (np7) -- (cpu);
    \draw[<->, thick] (cp9) -- (cpu);
    \draw[<->, thick] (mem) -- (cpu);
    \draw[<->, thick] (net1) -- (np7);
    \draw[<->, thick] (net2) -- (cp9);
\end{tikzpicture}
\caption{Architecture simplifiée d'un FortiGate --- Processeurs spécialisés}
\label{fig:architecture-fgt}
\end{figure}

\subsubsection{NP7 (Network Processor)}
Le NP7 est un processeur réseau dédié au traitement des paquets à haute vitesse \cite{1}. Il gère :
\begin{itemize}
    \item Le routage de base (forwarding)
    \item Le NAT (Network Address Translation)
    \item Le délestage de charge (offloading) pour IPSec
    \item La gestion des sessions réseau
\end{itemize}

\subsubsection{CP9 (Content Processor)}
Le CP9 est un processeur de contenu spécialisé dans l'analyse approfondie \cite{1}:
\begin{itemize}
    \item L'inspection antivirus et antispam
    \item L'analyse IPS (Intrusion Prevention System)
    \item Web and DNS filtering
    \item L'inspection SSL/TLS
\end{itemize}

\subsection{Cycle de traitement des paquets}

Le traitement des paquets dans FortiOS suit un flux optimisé :

\begin{figure}[H]
\centering
\begin{tikzpicture}[
    node distance=0.8cm,
    flowstep/.style={rectangle, draw, fill=green!10, text width=2.2cm, minimum height=0.7cm, align=center, font=\footnotesize},
    arrow/.style={-{Stealth[length=3mm]}, thick},
    scale=0.85,
    transform shape
]
    \node[flowstep] (input) {Entree};
    \node[flowstep, right=of input] (np) {NP7\\Routage};
    \node[flowstep, right=of np] (policy) {Politiques\\Firewall};
    \node[flowstep, right=of policy] (cp) {CP9\\UTM};
    \node[flowstep, right=of cp] (output) {Sortie};
    
    \draw[arrow] (input) -- (np);
    \draw[arrow] (np) -- (policy);
    \draw[arrow] (policy) -- (cp);
    \draw[arrow] (cp) -- (output);
\end{tikzpicture}
\caption{Cycle de traitement des paquets dans FortiOS}
\label{fig:paquet-lifecycle}
\end{figure}

\subsection{Flow-based vs Proxy-based}

FortiOS supporte deux modes d'inspection :

\begin{table}[H]
\centering
\caption{Comparaison des modes d'inspection}
\label{tab:flow-vs-proxy}
\begin{tabularx}{\textwidth}{|l|X|X|}
\hline
\textbf{Caractéristique} & \textbf{Flow-based} & \textbf{Proxy-based} \\
\hline
Performance & Élevée & Moyenne \\
\hline
Latence & Faible & Plus élevée \\
\hline
Fonctionnalités & UTM de base & UTM avancé \\
\hline
Utilisation & Routage, VPN & Inspection SSL, antivirus \\
\hline
\end{tabularx}
\end{table}

\section{Protocoles de routage}

\subsection{Routage statique}

Le routage statique consiste à définir manuellement les chemins empruntés par le trafic. Chaque route est caractérisée par :
\begin{itemize}
    \item \textbf{Destination} : sous-réseau cible
    \item \textbf{Passerelle} (gateway) : prochain saut
    \item \textbf{Distance administrative} (AD) : fiabilité de la source de routage
    \item \textbf{Métrique} : coût du chemin
\end{itemize}

Dans notre laboratoire, chaque FortiGate dispose de 3 routes maximum (contrainte de licence) :
\begin{enumerate}
    \item Route par défaut via WAN
    \item Route vers le sous-réseau LAN local
    \item Route vers le sous-réseau transit
\end{enumerate}

\subsection{OSPF (Open Shortest Path First)}

OSPF est un protocole de routage à état de liens (link-state) qui :
\begin{itemize}
    \item Échange des informations d'état de liens entre neighboring routers
    \item Construit une base de données topologique (LSDB)
    \item Calcule le plus court chemin (algorithme de Dijkstra)
    \item Réagit rapidement aux changements de topologie
\end{itemize}

Dans une topologie de type site-à-site, OSPF permettrait le routage entre les deux FortiGate via un lien transit (10.0.0.0/30) \cite{10}.

\subsection{BGP (Border Gateway Protocol)}

BGP est le protocole de routage inter-domaine utilisé sur Internet. Il échange des informations de routage entre systèmes autonomes (AS) \cite{10}.

Dans notre laboratoire, BGP n'a pas pu être pleinement démontré en raison de la limite de 3 routes imposée par la licence d'évaluation. Seule une vue conceptuelle du protocole est présentée dans cette section.

\section{High Availability (HA)}

\subsection{Protocole FGCP}

Le protocole FGCP (FortiGate Clustering Protocol) gère la formation et la synchronisation des clusters FortiGate \cite{2}. Il permet :
\begin{itemize}
    \item La détection automatique des pairs
    \item L'élection du noeud maître (master)
    \item La synchronisation des sessions et de la configuration
    \item Le failover automatique en cas de défaillance
\end{itemize}

\subsection{Active-Passive vs Active-Active}

\begin{table}[H]
\centering
\caption{Comparaison des modes HA}
\label{tab:ha-modes}
\begin{tabularx}{\textwidth}{|l|X|X|}
\hline
\textbf{Mode} & \textbf{Active-Passive} & \textbf{Active-Active} \\
\hline
Fonctionnement & Un noeud actif, un en attente & Les deux noeuds actifs \\
\hline
Charge & Uniquement sur le master & Distribuée \\
\hline
Synchronisation & Config + sessions & Config uniquement \\
\hline
Failover & Automatique & N/A \\
\hline
Utilisation & La plus courante & Charge élevée \\
\hline
\end{tabularx}
\end{table}

\section{VPN}

\subsection{IPSec VPN}

IPSec (Internet Protocol Security) fournit une sécurisation des communications IP via \cite{3}:
\begin{itemize}
    \item \textbf{IKEv1/IKEv2} : échange de clés et établissement du tunnel (Phase 1)
    \item \textbf{ESP} (Encapsulating Security Payload) : chiffrement des données (Phase 2)
    \item \textbf{Tunnel interface} : interface virtuelle pour le routage
\end{itemize}

\subsection{SSL VPN}

Le SSL VPN utilise le protocole SSL/TLS pour sécuriser les connexions distantes \cite{4}:
\begin{itemize}
    \item \textbf{Portail Web} : interface d'accès via navigateur
    \item \textbf{Client FortiClient} : application dédiée
    \item \textbf{Interface ssl.root} : tunnel virtuel pour le routage
\end{itemize}

\section{Modules UTM}

\subsection{Antivirus}

L'antivirus FortiGate détecte et bloque les fichiers malveillants. En mode évaluation, les signatures factory permettent la détection de l'EICAR (test standard).

\subsection{IPS (Intrusion Prevention System)}

L'analyse IPS détecte les attaques réseau et applicatives : scan de ports, SQL injection, XSS. Les signatures factory couvrent les attaques courantes.

\subsection{Web Filtering}

Web filtering bloque l'accès à des URLs ou catégories spécifiques. Sans licence FortiGuard, les listes d'URL statiques sont utilisées.

\subsection{DNS Filtering}

DNS filtering bloque les domaines malveillants. En mode statique, les domaines sont ajoutés manuellement à une liste de blocage.

\subsection{SSL Inspection}

L'inspection SSL/TLS déchiffre le trafic HTTPS pour analyse. Elle utilise un certificat CA auto-signé pour établir la confiance \cite{5}.
